\section*{ВВЕДЕНИЕ}
\addcontentsline{toc}{section}{ВВЕДЕНИЕ}

Современное развитие транспортной инфраструктуры требует постоянного контроля состояния дорожного покрытия и своевременного выявления возникающих дефектов. Рост протяжённости автомобильных дорог и интенсивности движения приводит к увеличению объёмов диагностических работ, что делает задачу автоматизации анализа дорожного полотна особенно актуальной.

Традиционные методы обследования дорог в значительной степени основаны на ручном анализе и визуальном контроле. Подобные подходы требуют существенных временных затрат и обладают ограниченной эффективностью при обработке больших объёмов данных. В связи с этим всё более широкое распространение получают программно-информационные системы, использующие методы компьютерного зрения и машинного обучения для автоматического анализа изображений.

Современные нейросетевые модели обнаружения объектов позволяют выполнять обработку изображений с высокой скоростью и точностью. Использование таких технологий делает возможным создание прикладных программных систем, способных автоматически выявлять дефекты дорожного покрытия на цифровых изображениях и формировать результаты анализа в удобной для пользователя форме.

Разработка подобных систем представляет практический интерес как для организаций, осуществляющих мониторинг дорожной инфраструктуры, так и для исследовательских задач в области компьютерного зрения.

\emph{Цель настоящей работы} -- создание программно-информационной системы обнаружения дефектов дорожного покрытия на основе нейросетевых методов анализа изображений. Для достижения поставленной цели необходимо решить \emph{следующие задачи}:

\begin{itemize}
	\item анализ предметной области и существующих методов обнаружения дефектов дорожного покрытия;
	\item изучение особенностей архитектур нейросетевых моделей обнаружения объектов;
	\item проектирование программной системы;
	\item обучение нейросетевой модели;
	\item проектирование пользовательского интерфейса приложения;
	\item реализация обработки изображений с использованием нейросетевой модели;
	\item реализация механизмов сохранения результатов и формирования отчёта.
\end{itemize}

\emph{Структура и объем работы.} Отчет состоит из введения, 4 разделов основной части, заключения, списка использованных источников, 2 приложений. Текст выпускной квалификационной работы равен \formbytotal{lastpage}{страниц}{е}{ам}{ам}.

\emph{Во введении} обоснована актуальность темы выпускной квалификационной работы, сформулированы цель и задачи разработки, приведена структура работы.

\emph{В первом разделе} выполнен анализ предметной области, рассмотрены существующие подходы к обнаружению дефектов дорожного покрытия, методы компьютерного зрения и архитектуры нейросетевых моделей.

\emph{Во втором разделе} представлено техническое задание на разработку программно-информационной системы, сформулированы требования к функциональным возможностям, интерфейсу.

\emph{В третьем разделе} описан технический проект системы, приведено обоснование выбора технологий разработки, рассмотрена архитектура программной системы и особенности внедрения нейросетевой модели.

\emph{В четвертом разделе} представлено описание программной реализации, приведены сведения о тестировании и проверке работоспособности разработанной системы.

В заключении приведены основные результаты работы и выводы по итогам разработки программно-информационной системы обнаружения дефектов дорожного покрытия.

В приложении А представлен графический материал.
В приложении Б представлены фрагменты исходного кода программы.


